\chapter{Çoktan Seçmeli Soruların Cevap Anahtarları}

Çoktan seçmeli soru bulunmayan Bölüm~\ref{chap:sentez} listede yer almamaktadır.

\begingroup
\small
\newcommand{\answerkey}[2]{\par\medskip\noindent\textbf{#1}\par\noindent #2\par}

\answerkey{Bölüm~\ref{chap:matematigin-dogasi} — Matematiğin Doğası ve Matematiksel Düşünme}{1-B, 2-B, 3-B, 4-C, 5-B, 6-B, 7-B, 8-B, 9-A, 10-B, 11-A, 12-B, 13-A, 14-B, 15-B, 16-B, 17-B, 18-C, 19-B, 20-B.}






\answerkey{Bölüm~\ref{chap:muhakeme-turleri} — Muhakeme Türleri}{1-B, 2-B, 3-C, 4-B, 5-B, 6-B, 7-B, 8-B, 9-B, 10-C.}

\answerkey{Bölüm~\ref{chap:muhakeme-ogretimi} — Düşünme Düzeyleri ve Muhakeme Öğretimi}{1-C, 2-A, 3-C, 4-C, 5-B, 6-B, 7-C, 8-A, 9-B, 10-B, 11-C, 12-A, 13-B.}



\answerkey{Bölüm~\ref{chap:onermeler-mantigi} — Önermeler Mantığı ve Argüman Geçerliliği}{1-C, 2-B, 3-A, 4-C, 5-B, 6-C, 7-C, 8-B, 9-B, 10-B, 11-C, 12-A, 13-B, 14-B, 15-C, 16-A, 17-B, 18-C, 19-B, 20-B.}

\answerkey{Bölüm~\ref{chap:sembolik-mantik} — Değişkenler}{1-B, 2-B, 3-B, 4-C, 5-B, 6-A, 7-B, 8-C, 9-B, 10-B, 11-B, 12-B, 13-B, 14-A, 15-C, 16-A, 17-B, 18-A, 19-A, 20-A.}

\answerkey{Bölüm~\ref{chap:niceleyiciler} — Niceleyiciler}{1-B, 2-B, 3-B, 4-B, 5-C, 6-A, 7-C, 8-B, 9-A, 10-B, 11-C, 12-B, 13-A, 14-C, 15-A, 16-A, 17-B, 18-C, 19-A, 20-B.}

\answerkey{Bölüm~\ref{chap:dogrudan-ispat} — İspat Kavramı; Doğrudan ve Kontrapozitif İspat}{1-B, 2-B, 3-B, 4-B, 5-B, 6-A, 7-B, 8-B, 9-B, 10-B, 11-A, 12-B, 13-C, 14-B, 15-B, 16-A, 17-B, 18-A, 19-A, 20-B.}

\answerkey{Bölüm~\ref{chap:celiski-ispat} — Çelişki ile İspat ve Niceleyicili İspatlar}{1-B, 2-D, 3-A, 4-C, 5-D, 6-A, 7-B, 8-C, 9-D, 10-A, 11-B, 12-C, 13-D, 14-A, 15-B, 16-C, 17-D, 18-A, 19-B, 20-C.}

\answerkey{Bölüm~\ref{chap:tumevarim} — Tümevarımla İspat}{1-B, 2-B, 3-C, 4-B, 5-B, 6-C, 7-B, 8-B, 9-B, 10-A, 11-B, 12-B, 13-B, 14-C, 15-B, 16-B, 17-B, 18-B, 19-B, 20-B.}

\answerkey{Bölüm~\ref{chap:kumeler} — Kümeler ve Küme Üzerine Muhakeme}{1-B, 2-B, 3-D, 4-B, 5-C, 6-B, 7-C, 8-B, 9-C, 10-A, 11-B, 12-B, 13-C, 14-C, 15-A, 16-B, 17-A, 18-B, 19-A, 20-A.}

\answerkey{Bölüm~\ref{chap:bagintilar} — Bağıntılar ve Denklik}{1-B, 2-B, 3-B, 4-B, 5-A, 6-B, 7-C, 8-D, 9-B, 10-B, 11-B, 12-A, 13-B, 14-B, 15-B, 16-B, 17-B, 18-B, 19-A, 20-B.}

\answerkey{Bölüm~\ref{chap:kongruanslar} — Kongrüanslar}{1-B, 2-B, 3-B, 4-C, 5-A, 6-C, 7-D, 8-B, 9-C, 10-B, 11-C, 12-B, 13-B, 14-B, 15-B, 16-B, 17-B, 18-C, 19-B, 20-B.}

\answerkey{Bölüm~\ref{chap:fonksiyonlar} — Fonksiyonlar}{1-B, 2-C, 3-A, 4-B, 5-B, 6-C, 7-D, 8-A, 9-B, 10-B, 11-C, 12-B, 13-A, 14-B, 15-B, 16-C, 17-B, 18-B, 19-C, 20-B.}

\answerkey{Bölüm~\ref{chap:kardinalite} — Kardinalite ve Sonsuzluk}{1-B, 2-C, 3-B, 4-B, 5-B, 6-A, 7-A, 8-B, 9-B, 10-B, 11-B, 12-B, 13-B, 14-B, 15-B, 16-C, 17-B, 18-B, 19-B, 20-B.}



\endgroup
